%% ===========================================================================
\section{Conclusion}
\budget{250}

\jun{Prose draft (my part; old bullet spec commented below). Follows the spec:
thesis in the two strongest number-pairs, one paragraph on what the negatives
buy, one concrete next step and no laundry list. Dropped the labeller-comparison
item from the ``negatives buy'' sentence to stay consistent with the abstract,
where that contribution is cut (Mohit). Numbers checked against \S\ref{sec:results}.
Yash: two sentences characterise your results (R5 positivity 0.094 vs 0.050 in
para 1; R6 bridge-user continuum in para 2) --- the figures match
\S\ref{sec:results} but please confirm the wording.}

Across six India--Pakistan crises, what a community contested turned on the frame
a comment advanced and the community it reached, not on how offensively it was
worded: humanitarian framing split a general-audience community 17.4\% of the
time, where offensive and inoffensive comments diverged by only 7.67\% against
6.83\%. The same frame was not the same object across communities ---
security-and-military framing read as positive at 0.094 in India-aligned
communities against 0.050 in general-audience ones --- so community governs
\emph{what} a discourse contests and frame governs \emph{how} it is voiced, and
neither axis is visible to a comment scored on a single layer.

Two of these results exist only because the paper reported its failures. The
difference-in-differences designs collapsed under a wild cluster bootstrap, which
turned a null into a transferable caution: crisis-response causal inference on a
handful of temporally clustered events is uninformative, not merely underpowered.
And the search for a bridge-user type returned a continuum rather than a kind,
which is why we describe boundary-crossing as a behaviour with a distribution
rather than a class of person.

One next step follows directly. The veracity of the disputed factual claims that
drive the most contested cell is the layer we most want and do not yet have;
adding it means first repairing a relevance gate that currently recovers only
4.8\% of those claims.

% OLD BULLET SPEC (superseded by prose above):
% \begin{itemize}
% \item Restate the thesis with the two numbers that do the most work: neutral
%   x Humanitarian 17.4\% against the Offensive-versus-not gap of 7.67\% / 6.83\%;
%   and the Security/Military positivity asymmetry (0.094 india versus 0.050 neutral).
% \item One sentence on what the negative results buy --- the event-clustering
%   confounder, the continuum finding, and the labeller comparison are all results
%   the paper only has because it reported failures.
% \item No future-work laundry list. One concrete next step: the L7 veracity layer
%   over the gated Factual Claims, which requires first fixing an L1 gate that
%   currently recovers 4.8\% of them.
% \end{itemize}
